\section{Optimal integral frames and the ramification obstruction}
\label{sec:frames}

For a totally real number field $F$ of degree $n$, an ordinary integral tight frame of bound $M\in\Z_{>0}$ is a finite family $v_i\in\OO_F$ with positive integer multiplicities $m_i$ such that
\begin{equation}\label{eq:frame}
 \sum_i m_i\sigma(v_i)\tau(v_i)=M\delta_{\sigma,\tau}
 \qquad(\sigma,\tau:F\hookrightarrow\R).
\end{equation}
Thus the vectors give a resolution of a scalar identity in the embedding space. Repeated vectors can replace multiplicities. This is the linear-algebra part of the source's frame input; its global integrality condition instead uses a fractional Weil lattice.

\subsection{Why a ramified prime must divide the bound}

\begin{theorem}[Trace-dual obstruction]\label{thm:dual-obstruction}
If \eqref{eq:frame} holds, then
\begin{equation}\label{eq:codiff-obstruction}
 M\DD_{F/\Q}^{-1}\subseteq\OO_F.
\end{equation}
In particular, every rational prime ramified in $F$ divides $M$.
\end{theorem}

\begin{proof}
Let $a$ lie in the codifferent, so $\Tr_{F/\Q}(a v_i)\in\Z$ for all $i$. Applying the frame operator to the embedding vector of $a$ gives
\[
 M a=\sum_i m_i\Tr_{F/\Q}(a v_i)v_i\in\OO_F.
\]
This proves \eqref{eq:codiff-obstruction}. At a prime $\mathfrak p$ over a ramified rational prime $p$, write $e_{\mathfrak p}$ for the ramification index and $d_{\mathfrak p}>0$ for the different exponent. The ideal containment requires
\[
 e_{\mathfrak p}v_p(M)\ge d_{\mathfrak p},
\]
so $v_p(M)\ge1$.
\end{proof}

The obstruction concerns every ordinary integral frame, not only a particular construction. It also explains why the source allows denominators in one half of the relative codifferent. Those denominators change the dual lattice; they are essential arithmetic data.

\subsection{The exact eventual congruence obstruction for an integral lattice}

The elementary obstruction has a useful converse away from a finite initial range.

\begin{theorem}[Eventual frame bounds for an integral lattice]\label{thm:real-eventual}
Let $L$ be a positive definite integral $\Z$-lattice of positive rank with Gram matrix $G$ in an integral basis. Put
\[
 d=\min\{a\in\Z_{>0}:aG^{-1}\text{ has integer entries}\}.
\]
Every positive integer tight-frame bound for vectors in $L$ is divisible by $d$, and every sufficiently large positive multiple of $d$ is such a bound.
\end{theorem}

\begin{proof}
Let $B$ be the real basis matrix, so $G=B^{\mathsf T}B$. If the frame vectors have integer coordinate columns $u_i$, its identity becomes
\[
 \sum_i m_i u_i u_i^{\mathsf T}=M G^{-1}.
\]
The left side is integral. This proves $d\mid M$; equivalently, $d$ is the exponent of the discriminant group $L^\#/L$.

For sufficiency set $C=dG^{-1}$, a positive definite integral symmetric matrix. The rank-one matrices generated by $e_i$ and $e_i+e_j$ span all integral symmetric matrices over $\Z$: their differences give the off-diagonal matrix units, while $e_ie_i^{\mathsf T}$ gives the diagonal units. Let $S$ be this finite set of vectors. For a sufficiently small positive rational $\epsilon$, the matrix
\[
 C-\epsilon\sum_{v\in S}vv^{\mathsf T}
\]
is positive definite and rational. Rational $LDL^{\mathsf T}$ decomposition expresses it as a sum of positive rational multiples of rational rank-one matrices. Scaling their vectors to integer vectors and clearing coefficients yields
\[
 qC=\sum_{v\in T}c_vvv^{\mathsf T},\qquad c_v\in\Z_{>0},\quad S\subseteq T.
\]
Write $C=\sum_{v\in T}z_vvv^{\mathsf T}$ with $z_v\in\Z$, using $S$, and extending by zero on other vectors. For $n=kq+r$, $0\le r<q$,
\[
 nC=\sum_{v\in T}(kc_v+rz_v)vv^{\mathsf T}.
\]
For all sufficiently large $k$, every coefficient is nonnegative. The resulting frame has bound $nd$. The same argument makes every coefficient exceed any prescribed fixed lower bound.
\end{proof}

This proof supplies an effective threshold once a rational reserve decomposition has been chosen. It supplies no useful uniform bound as the rank and arithmetic data vary. Sums of integral squares and Hermitian norms form an established subject; the present formulation is used as a transfer tool.\cite{BCIL}

\subsection{The cubic field and its optimal six-vector certificate}

Let $f(X)=X^3-4X+1$ and let $\alpha$ be a root. The rational root test proves irreducibility. There is a root in each of $(-3,-2)$, $(0,1)$, and $(1,2)$, so the cubic field $F_0=\Q(\alpha)$ is totally real. Its polynomial discriminant is
\[
 -4(-4)^3-27=229,
\]
a prime. Thus $(1,\alpha,\alpha^2)$ is an integral basis, and the trace Gram matrix is
\[
 T=\begin{pmatrix}3&0&8\\0&8&-3\\8&-3&32\end{pmatrix},
 \qquad
 229T^{-1}=\begin{pmatrix}247&-24&-64\\-24&32&9\\-64&9&24\end{pmatrix}.
\]
Newton sums give $\Tr(1),\Tr(\alpha),\ldots,\Tr(\alpha^4)=(3,0,8,-3,32)$.

The following coordinate vectors have outer products summing to $229T^{-1}$:
\begin{center}
\begin{tabular}{@{}rrrcr@{}}
\toprule
$1$&$\alpha$&$\alpha^2$&&Multiplicity\\\midrule
0&2&1&&1\\
1&-1&1&&2\\
2&-1&1&&1\\
4&-5&-2&&1\\
15&0&-4&&1\\\bottomrule
\end{tabular}
\end{center}
Therefore they give a six-vector frame of bound $229$. Theorem~\ref{thm:dual-obstruction} shows that no smaller positive integer bound is possible. No assertion about a minimum number of vectors is needed.

\subsection{The $S_3$ field and an optimal eight-vector certificate}

Let $F$ be the splitting field and let $\alpha,\beta$ be distinct roots of $f$. Its Galois group is $S_3$, because the irreducible cubic has nonsquare discriminant. Its six embeddings correspond to ordered pairs of distinct roots. Define
\[
 b=(1,\alpha,\alpha^2,\beta,\alpha\beta,\alpha^2\beta).
\]
This is an integral basis. Here is a direct discriminant justification, also useful below. Only $229$ can ramify in the splitting field. At $229$ the extension is tame, and
\[
 f(X)\equiv(X-29)^2(X+58)\pmod{229}.
\]
Inertia has order two and acts as a transposition: it is nontrivial, cyclic, and fixes the simple-root branch. A tame degree-six Galois extension with inertia index two has discriminant exponent $6(1-1/2)=3$. Consequently $\operatorname{disc}(F)=229^3$. The Gram matrix of $b$, computed below, also has determinant $229^3$, proving that the integral lattice it generates has index one in $\OO_F$.

Write $s_j$ for the $j$th power sum of the three roots. Then
\[
 s_0=3,\quad s_1=0,\quad s_2=8,\qquad
 s_j=4s_{j-2}-s_{j-3}\quad(j\ge3).
\]
For monomials $\alpha^i\beta^j$ and $\alpha^k\beta^l$, their trace product equals
\[
 s_{i+k}s_{j+l}-s_{i+j+k+l}.
\]
It follows that the Gram matrix and its scaled inverse are
\begin{align*}
G&=\begin{pmatrix}
6&0&16&0&-8&3\\
0&16&-6&-8&3&-32\\
16&-6&64&3&-32&20\\
0&-8&3&16&3&32\\
-8&3&-32&3&32&-4\\
3&-32&20&32&-4&125
\end{pmatrix},\\[3pt]
C:=229G^{-1}&=\begin{pmatrix}
125&-16&-32&4&0&-3\\
-16&32&3&0&-3&8\\
-32&3&16&-3&8&0\\
4&0&-3&32&-6&-8\\
0&-3&8&-6&16&0\\
-3&8&0&-8&0&6
\end{pmatrix}.
\end{align*}

\begin{theorem}[Optimal discriminant-229 frame]\label{thm:optimal229}
Let $v_i$ have the following coordinates in $b$, with multiplicities $m_i$:
\begin{center}
\begin{tabular}{@{}rrrrrrcr@{}}
\toprule
$1$&$\alpha$&$\alpha^2$&$\beta$&$\alpha\beta$&$\alpha^2\beta$&&$m_i$\\\midrule
0&2&-2&-2&-1&1&&1\\
1&-2&-2&4&-2&-1&&1\\
2&2&0&1&2&1&&1\\
4&1&-1&1&-1&0&&2\\
4&1&-1&2&0&-1&&1\\
6&-4&-1&-1&2&-1&&1\\
6&-1&-2&-2&-1&1&&1\\\bottomrule
\end{tabular}
\end{center}
They satisfy \eqref{eq:frame} with $M=229$, and $229$ is the smallest possible positive integer bound for any ordinary integral frame in $F$.
\end{theorem}

\begin{proof}
Let $u_i$ denote the displayed columns. Direct integer multiplication gives
\[
 \sum_i m_i u_i u_i^{\mathsf T}=C,
 \qquad GC=229I_6.
\]
If $B=(\sigma(b_j))_{\sigma,j}$, then $B^{\mathsf T}B=G$, whence
\[
 \sum_i m_i\sigma(v_i)\tau(v_i)
 =(BCB^{\mathsf T})_{\sigma,\tau}=229\delta_{\sigma,\tau}.
\]
The vectors are integral. Every positive integer frame bound is divisible by the ramified prime $229$, by Theorem~\ref{thm:dual-obstruction}. This proves optimality.
\end{proof}

The squared trace norms are, in order,
\[
 (221,229,229,74,157,211,179),
\]
and their multiplicity-weighted sum is $6\cdot229$. In particular all real conjugates of every displayed vector have absolute value at most $\sqrt{229}$. The exact checker recomputes the Newton sums, matrices, determinants, outer products, and identities. A numerical search was used only to discover candidate integer columns; no numerical optimizer is part of the verification.

\begin{corollary}[Cyclotomic tensor frame]
For $229\mid q$, the elements $v_i\zeta_q^j$, $0\le j<q$, with the same multiplicities, give a frame of bound $229q$ on the compatible embeddings of $F\Q(\zeta_q)$. The same embedding identity holds on the induced CM type after adjoining the source's imaginary quadratic field $K$.
\end{corollary}

\begin{proof}
Multiply \eqref{eq:frame} by the geometric-sum identity
\[
 \sum_{j=0}^{q-1}\eta(\zeta_q^j)\overline{\eta'(\zeta_q^j)}
 =q\delta_{\eta,\eta'}.
\]
Embeddings of the compositum correspond to compatible pairs of embeddings. Restricting the product identity to those pairs proves the assertion; linear disjointness is not assumed.
\end{proof}

The tensor frame improves an explicit arithmetic constant and preserves integrality at the inert primes used in the earlier local argument. Its bound still carries $229$. The correction needed for coprime bounds is a change to the appropriate fractional lattice, addressed next.
